\subsection*{Процедура Merge}

\Task Даны два отсортированных массива $a_1 \leq\,\dots\,\leq a_n$ и
$b_1 \leq\,\dots\,\leq b_m$. Надо слить их в один отсортированный за O(n + m)
времени и доппамяти.

\Solve 
\begin{enumerate}
    \item Заведем указатели на первые элементы массивов i и j.
    \item По индексу i + j выпишем меньший из $a_i$ и $b_j$ и сдвинем соответствующий указатель на один вперед.
    \item Так делаем, пока i < n \&\& j < m.
    \item Оставшийся хвост одного из массивов дописываем в конец.
\end{enumerate}

\subsection*{MergeSort}

\Task Дан массив $a_1, \dots , a_n$. Необходимо отсортировать его за
$O(n\,log\,n)$ времени и $O(n)$ доппамяти

\Solve
\begin{enumerate}
    \item Разделим массив на две равные части.
    \item Каждую из них рекурсивно отсортируем.
    \item Результаты сольем в один отсортированный массив с помощью процедуры Merge.
\end{enumerate}

\subsection*{Инверсии}

\Def Пара i < j называется инверсией, если $a_i > a_j$.

\Task Дан массив $a_1, \dots , a_n$. Необходимо найти число инверсий за
$O(n\,log\,n)$ времени и $O(n)$ доппамяти.

\Solve 
\begin{enumerate}
    \item Если при слиянии элемент левой части больше элемента правой части, то значит это и есть инверсия.
    \item Пусть мы сравниваем в сортировке слиянием l[i] и r[j], тогда если $r[j] < l[i]$, то $r[j] < l[i] < l[i + 1] <\, \dots \,< l[N]$, то есть число инверсий это N - i для r[j].
    \item Тогда для одного слияния итоговое число инверсий равно сумме по j таких значений.
    \item Итоговый ответ равен сумме по всем слияниям.
\end{enumerate}

\pagebreak